\documentclass[a4paper,12pt]{article}

\usepackage[T2A]{fontenc}
\usepackage[utf8]{inputenc}
\usepackage[russian]{babel}
\usepackage{amsmath,amssymb}
\usepackage{PTSans}
\renewcommand{\familydefault}{\sfdefault}
\usepackage{icomma}
\usepackage[a4paper,margin=18mm,headheight=15pt,footskip=12mm]{geometry}
\usepackage{microtype}
\usepackage{tikz}
\usetikzlibrary{calc,arrows.meta,positioning,shapes.geometric}
\usepackage{tkz-euclide}
\usepackage{tabularx,booktabs}
\usepackage{enumitem}
\usepackage{fancyhdr}
\usepackage{setspace}
\usepackage{xcolor}

\definecolor{accent}{HTML}{008080}
\definecolor{darkaccent}{HTML}{004D4D}
\definecolor{textgray}{HTML}{333333}

\color{textgray}
\renewcommand{\today}{26 сентября 2026 г.}
\onehalfspacing
\setlength{\parindent}{0pt}
\setlength{\parskip}{0.55em}
\renewcommand{\arraystretch}{1.28}

\setlist[itemize]{
    leftmargin=1.6em,
    itemsep=0.25em,
    topsep=0.35em
}

\setlist[enumerate]{
    leftmargin=2em,
    itemsep=0.35em,
    topsep=0.35em
}

\pagestyle{fancy}
\fancyhf{}
\fancyhead[L]{\small Екатерина Скелли}
\fancyhead[R]{\small Подготовка к МЦКО}
\fancyfoot[C]{\small\thepage}
\renewcommand{\headrulewidth}{0.3pt}
\renewcommand{\footrulewidth}{0pt}

\newcommand{\sectionline}{%
    \vspace{-0.15em}
    {\color{accent}\hrule height 0.7pt}
    \vspace{0.55em}
}

\newcommand{\keyidea}[1]{%
    \par\noindent
    {\color{darkaccent}\textbf{Ключевая идея.}} #1\par
}

\newcommand{\warning}[1]{%
    \par\noindent
    {\color{darkaccent}\textbf{Проверьте себя.}} #1\par
}

\newcommand{\formula}[1]{%
    \[
        \boxed{\displaystyle #1}
    \]
}

\tkzSetUpLine[line width=0.9pt,color=accent]
\tkzSetUpPoint[color=accent,fill=accent,size=3]
\tkzSetUpLabel[color=black,font=\small]

\tikzset{
    fcstart/.style={
        ellipse,
        draw=darkaccent,
        line width=0.9pt,
        minimum width=4.2cm,
        minimum height=1cm,
        align=center,
        font=\small
    },
    fcaction/.style={
        rectangle,
        rounded corners=3pt,
        draw=accent,
        line width=0.9pt,
        text width=8.3cm,
        minimum height=1.05cm,
        align=center,
        font=\small
    },
    fcdecision/.style={
        diamond,
        aspect=2.55,
        draw=accent,
        line width=0.9pt,
        text width=5.4cm,
        align=center,
        inner sep=1.5pt,
        font=\small
    },
    fcarrow/.style={
        -{Stealth[length=2.8mm,width=2mm]},
        line width=0.9pt,
        draw=textgray
    }
}

\begin{document}

\begin{titlepage}
\thispagestyle{empty}

\begin{center}
\vspace*{16mm}

{\color{darkaccent}\Large\bfseries ЧЕК-ЛИСТ}

\vspace{7mm}

{\Huge\bfseries
Комплексное повторение\\[2mm]
перед МЦКО}

\vspace{6mm}

{\Large Геометрия и статистика}

\vspace{18mm}

{\large Екатерина Скелли}

\vspace{12mm}

{\color{accent}\rule{0.54\linewidth}{0.9pt}}

\vspace{12mm}

\begin{minipage}{0.80\linewidth}
\centering
\large
Треугольники и углы $\bullet$
параллельные прямые $\bullet$
прямоугольный треугольник $\bullet$
окружность $\bullet$
неравенство треугольника $\bullet$
среднее, медиана и мода $\bullet$
круговые диаграммы
\end{minipage}

\vfill

{\large \today}

\vspace{10mm}

{\small
Лёвин Артём Александрович\\
эксклюзивно для Екатерины Скелли
}

\vspace{16mm}

\begin{tikzpicture}[remember picture,overlay]
\draw[
    accent,
    line width=1.05pt
]
([xshift=18mm,yshift=17mm]current page.south west)
.. controls
([xshift=56mm,yshift=26mm]current page.south west)
and
([xshift=-56mm,yshift=8mm]current page.south east)
..
([xshift=-18mm,yshift=17mm]current page.south east);
\end{tikzpicture}

\end{center}
\end{titlepage}

\section*{\color{darkaccent}1. Как начинать геометрическую задачу}
\sectionline

На занятии основной акцент был сделан на геометрии. При повторении полезно
придерживаться одной и той же последовательности рассуждений.

\begin{enumerate}
    \item Отметьте на чертеже то, что требуется найти.
    \item Найдите треугольник, в который искомый отрезок или угол входит как
    полноценный элемент.
    \item Определите специальный вид этого треугольника: равнобедренный,
    прямоугольный, треугольник с углом $30^\circ$ и т.~д.
    \item Подключите свойства параллельных прямых, биссектрис, медиан,
    высот, окружности, если они есть в условии.
    \item Если одна и та же величина выражается двумя способами, приравняйте
    полученные выражения и составьте уравнение.
\end{enumerate}

\keyidea{Чертёж помогает увидеть связи. Обоснование решения должно опираться
на теоремы и свойства, а не на внешний вид рисунка.}

\section*{\color{darkaccent}2. Неравенство треугольника}
\sectionline

Для сторон любого невырожденного треугольника выполняются три строгих
неравенства:

\[
a+b>c,\qquad a+c>b,\qquad b+c>a.
\]

Если две стороны равны $a$ и $b$, а третья равна $x$, все три условия удобно
свернуть в одно:

\formula{|a-b|<x<a+b.}

\subsection*{Пример}

Две стороны равны $7$ и $11$. Найдём все положительные целые значения третьей
стороны $x$:

\[
|11-7|<x<11+7,
\]

\[
4<x<18.
\]

Следовательно,

\[
x\in\{5,6,7,\ldots,17\}.
\]

Количество значений:

\[
17-5+1=13.
\]

\begin{center}
\begin{tikzpicture}[x=0.48cm,y=1cm]
    \draw[-{Stealth[length=2.5mm]},line width=0.8pt] (0,0)--(19,0);

    \foreach \x in {1,...,18}{
        \draw (\x,0.09)--(\x,-0.09);
    }

    \foreach \x in {1,2,3,4,5,10,15,17,18}{
        \node[below=4pt,font=\small] at (\x,0) {\x};
    }

    \draw[accent,line width=2pt] (5,0.28)--(17,0.28);
    \draw[accent,fill=white,line width=1pt] (4,0.28) circle (2.4pt);
    \draw[accent,fill=white,line width=1pt] (18,0.28) circle (2.4pt);

    \node[above,font=\small] at (11,0.36) {$4<x<18$};
\end{tikzpicture}
\end{center}

\warning{Границы не входят в ответ. При $x=|a-b|$ или $x=a+b$ треугольник
вырождается.}

\section*{\color{darkaccent}3. Равнобедренный треугольник}
\sectionline

Если $AB=BC$, то $AC$ --- основание равнобедренного треугольника $ABC$.

Нужно помнить:

\begin{itemize}
    \item углы при основании равны;
    \item медиана, проведённая из вершины к основанию, одновременно является
    высотой и биссектрисой;
    \item если периметр равен $P$, а основание --- $a$, то каждая боковая
    сторона равна
    \[
        \frac{P-a}{2}.
    \]
\end{itemize}

Например, при $P=32$ см и основании $10$ см боковая сторона равна

\[
\frac{32-10}{2}=11\text{ см}.
\]

\section*{\color{darkaccent}4. Как выразить один отрезок двумя способами}
\sectionline

Пусть $AB=BC$, $\angle ABC=120^\circ$, $BM$ --- медиана к основанию $AC$.
На луче $BM$ за точкой $M$ выбрана точка $F$, причём
$AF\perp AB$ и $MF=72$.

\begin{center}
\begin{tikzpicture}[scale=0.92]
    \tkzDefPoint(0,0){B}
    \tkzDefPoint(-2.598,-1.5){A}
    \tkzDefPoint(2.598,-1.5){C}
    \tkzDefPoint(0,-1.5){M}
    \tkzDefPoint(0,-6){F}

    \tkzDrawSegments(A,B B,C A,C B,F A,F)
    \tkzDrawPoints(A,B,C,M,F)

    \tkzLabelPoints[above](B)
    \tkzLabelPoints[left](A)
    \tkzLabelPoints[right](C)
    \tkzLabelPoints[right](M,F)

    \tkzMarkRightAngle[size=0.26](A,M,B)
    \tkzMarkRightAngle[size=0.26](B,A,F)

    \tkzMarkAngle[size=0.55](A,B,M)
    \tkzLabelAngle[pos=0.82](A,B,M){$60^\circ$}

    \tkzLabelSegment[left](A,B){$x$}
    \tkzLabelSegment[right](M,F){$72$}
\end{tikzpicture}
\end{center}

Обозначим $AB=x$. Так как $BM$ --- медиана к основанию равнобедренного
треугольника,

\[
BM\perp AC,
\qquad
\angle ABM=60^\circ.
\]

В прямоугольном треугольнике $ABM$

\[
\angle BAM=30^\circ,
\]

поэтому катет $BM$, лежащий против угла $30^\circ$, равен половине
гипотенузы:

\[
BM=\frac{x}{2}.
\]

Поскольку $F$ лежит на луче $BM$,

\[
\angle ABF=60^\circ.
\]

В прямоугольном треугольнике $ABF$ тогда

\[
\angle AFB=30^\circ.
\]

Сторона $AB$ лежит против угла $30^\circ$, следовательно,

\[
BF=2AB=2x.
\]

С другой стороны,

\[
BF=BM+MF=\frac{x}{2}+72.
\]

Приравниваем два выражения для $BF$:

\[
2x=\frac{x}{2}+72.
\]

Умножаем \emph{обе части уравнения} на $2$:

\[
4x=x+144,
\]

\[
3x=144,
\]

\[
\boxed{x=48}.
\]

\keyidea{Если одна и та же величина выражена двумя способами, равенство этих
выражений часто и есть нужное уравнение.}

\warning{При избавлении от знаменателя множитель применяется ко всей части
уравнения, то есть к каждому её слагаемому.}

\clearpage
\thispagestyle{empty}

\begin{center}
{\Large\bfseries\color{darkaccent}
Алгоритм: поиск неизвестного элемента в геометрической задаче}
\end{center}

\vspace{7mm}

\begin{center}
\begin{tikzpicture}[node distance=9mm]

\node[fcstart] (start) {Начало};

\node[fcaction,below=of start] (target)
{Отметьте искомый угол или отрезок};

\node[fcaction,below=of target] (triangle)
{Найдите треугольник, в который искомый элемент входит как сторона или угол};

\node[fcdecision,below=12mm of triangle] (special)
{Есть ли узнаваемое свойство или специальная конструкция?};

\node[fcaction,below=14mm of special,xshift=-40mm,text width=5cm] (apply)
{Примените подходящую теорему или свойство};

\node[fcaction,below=14mm of special,xshift=40mm,text width=5cm] (search)
{Найдите дополнительную связь: другой треугольник, смежный угол, параллельность};

\coordinate (joinone) at ($(apply.south)!0.5!(search.south)+(0,-22mm)$);

\node[fcdecision] (enough) at (joinone)
{Данных уже достаточно, чтобы найти искомое?};

\node[fcaction,below=14mm of enough,xshift=-40mm,text width=5cm] (calculate)
{Вычислите искомую величину};

\node[fcaction,below=14mm of enough,xshift=40mm,text width=5cm] (equation)
{Выразите одну величину двумя способами и составьте уравнение};

\coordinate (jointwo) at ($(calculate.south)!0.5!(equation.south)+(0,-16mm)$);

\node[fcaction,text width=8.3cm] (check) at (jointwo)
{Проверьте углы, длины, единицы и соответствие условию};

\node[fcstart,below=of check] (finish)
{Ответ};

\draw[fcarrow] (start)--(target);
\draw[fcarrow] (target)--(triangle);
\draw[fcarrow] (triangle)--(special);

\draw[fcarrow] (special)--node[above left,font=\small]{да}(apply);
\draw[fcarrow] (special)--node[above right,font=\small]{нет}(search);

\draw[fcarrow] (apply.south)|-(enough.west);
\draw[fcarrow] (search.south)|-(enough.east);

\draw[fcarrow] (enough)--node[above left,font=\small]{да}(calculate);
\draw[fcarrow] (enough)--node[above right,font=\small]{нет}(equation);

\draw[fcarrow] (calculate.south)|-(check.west);
\draw[fcarrow] (equation.south)|-(check.east);
\draw[fcarrow] (check)--(finish);

\end{tikzpicture}
\end{center}

\clearpage

\section*{\color{darkaccent}5. Смежные углы, отношения и биссектрисы}
\sectionline

Смежные углы в сумме дают

\formula{180^\circ.}

Если смежные углы относятся как $2:3$, обозначаем их $2x$ и $3x$:

\[
2x+3x=180^\circ,
\]

\[
5x=180^\circ,
\]

\[
x=36^\circ.
\]

Значит, сами углы равны

\[
72^\circ\quad\text{и}\quad108^\circ.
\]

\keyidea{Отношение $a:b$ удобно переводить в выражения $ax$ и $bx$. После
нахождения $x$ обязательно вернитесь к исходным углам.}

Биссектриса делит угол на две равные части. Если она образует с одной стороной
угла $27^\circ$, то весь угол равен

\[
27^\circ+27^\circ=54^\circ.
\]

Особенно полезный факт:

\formula{\text{биссектрисы двух смежных углов перпендикулярны}.}

Действительно, если смежные углы равны $\alpha$ и $\beta$, то

\[
\alpha+\beta=180^\circ,
\]

а угол между их биссектрисами равен

\[
\frac{\alpha}{2}+\frac{\beta}{2}
=
\frac{\alpha+\beta}{2}
=
90^\circ.
\]

\section*{\color{darkaccent}6. Параллельные прямые и секущая}
\sectionline

Если две прямые параллельны и пересечены секущей, работают три стандартные
связи.

\begin{center}
\begin{tabularx}{0.88\linewidth}{@{}X X@{}}
\toprule
\textbf{Тип углов} & \textbf{Свойство} \\
\midrule
Соответственные & равны \\
Накрест лежащие & равны \\
Односторонние & в сумме дают $180^\circ$ \\
\bottomrule
\end{tabularx}
\end{center}

\keyidea{Как только в условии появляется параллельность, сразу ищите
соответственные, накрест лежащие и односторонние углы.}

\section*{\color{darkaccent}7. Сумма углов и внешний угол треугольника}
\sectionline

Для любого треугольника

\formula{\angle A+\angle B+\angle C=180^\circ.}

Внешний угол треугольника равен сумме двух внутренних углов, не смежных с ним.
Например, если внешний угол при вершине $C$ равен $150^\circ$, а
$\angle B=120^\circ$, то

\[
150^\circ=\angle A+120^\circ,
\]

откуда

\[
\angle A=30^\circ.
\]

Внутренний угол при $C$ смежен с внешним, поэтому

\[
\angle C=180^\circ-150^\circ=30^\circ.
\]

Следовательно, в таком треугольнике

\[
\angle A=\angle C,
\]

а значит,

\[
AB=BC.
\]

\section*{\color{darkaccent}8. Прямоугольный треугольник}
\sectionline

В прямоугольном треугольнике два острых угла в сумме дают $90^\circ$.
Поэтому, если один острый угол равен $\alpha$, второй равен

\formula{90^\circ-\alpha.}

Например, при $\alpha=34^\circ$ второй угол равен

\[
90^\circ-34^\circ=56^\circ.
\]

\subsection*{Угол $30^\circ$}

Если в прямоугольном треугольнике острый угол равен $30^\circ$, то катет,
лежащий напротив этого угла, равен половине гипотенузы:

\formula{a=\frac{c}{2}.}

Верно и обратное утверждение: если катет равен половине гипотенузы, то угол,
лежащий напротив этого катета, равен $30^\circ$.

\subsection*{Угол $45^\circ$}

Если один острый угол равен $45^\circ$, то второй также равен $45^\circ$.
Следовательно, катеты равны.

\section*{\color{darkaccent}9. Медиана к гипотенузе}
\sectionline

В прямоугольном треугольнике медиана, проведённая из вершины прямого угла к
гипотенузе, равна половине гипотенузы.

\begin{center}
\begin{tikzpicture}[scale=1.05]
    \tkzDefPoint(0,0){B}
    \tkzDefPoint(4,0){A}
    \tkzDefPoint(0,3){C}
    \tkzDefMidPoint(A,C)
    \tkzGetPoint{M}

    \tkzDrawPolygon(A,B,C)
    \tkzDrawSegment(B,M)
    \tkzDrawPoints(A,B,C,M)

    \tkzLabelPoints[below](B,A)
    \tkzLabelPoints[left](C)
    \tkzLabelPoints[above right](M)

    \tkzMarkRightAngle[size=0.28](A,B,C)
\end{tikzpicture}
\end{center}

Если $M$ --- середина гипотенузы $AC$, то

\formula{BM=AM=CM=\frac{AC}{2}.}

Например, если $BM=6$ см, то

\[
AC=12\text{ см}.
\]

\warning{Это свойство относится именно к прямоугольному треугольнику и к
медиане, проведённой к гипотенузе.}

\section*{\color{darkaccent}10. Признаки равенства треугольников}
\sectionline

Для произвольных треугольников:

\begin{enumerate}
    \item по двум сторонам и углу между ними;
    \item по стороне и двум прилежащим к ней углам;
    \item по трём сторонам.
\end{enumerate}

Для прямоугольных треугольников часто удобно использовать сокращённые
признаки:

\begin{itemize}
    \item по двум катетам;
    \item по гипотенузе и катету;
    \item по катету и прилежащему к нему острому углу;
    \item по гипотенузе и острому углу.
\end{itemize}

Сокращённые признаки работают потому, что прямой угол уже известен и равен
$90^\circ$ в обоих треугольниках.

\section*{\color{darkaccent}11. Стороны, углы, высоты и расстояния}
\sectionline

В любом треугольнике:

\[
\text{против большей стороны лежит больший угол},
\]

\[
\text{против меньшей стороны лежит меньший угол}.
\]

Например, для сторон $7$, $9$ и $12$ наибольший угол лежит напротив стороны
$12$.

В тупоугольном треугольнике основание высоты может лежать на
\emph{продолжении} противоположной стороны. Высота определяется условием
перпендикулярности, поэтому она не обязана целиком находиться внутри
треугольника.

Расстояние от точки до прямой --- длина перпендикуляра, опущенного из точки на
эту прямую. Это кратчайший отрезок от точки до прямой.

Следовательно, если расстояние от точки $P$ до прямой равно $6$ см, наклонная
от $P$ к этой прямой не может иметь длину $5$ см.

\section*{\color{darkaccent}12. Окружность}
\sectionline

\subsection*{Радиус, диаметр и хорда}

\[
d=2r.
\]

Хорда, проходящая через центр окружности, является диаметром. Например, при
$r=7$ такая хорда имеет длину $14$.

\subsection*{Радиус и касательная}

Радиус, проведённый в точку касания, перпендикулярен касательной.

\begin{center}
\begin{tikzpicture}[scale=0.82]
    \tkzDefPoint(0,2){O}
    \tkzDefPoint(0,0){T}
    \tkzDefPoint(0,4){R}
    \tkzDefPoint(-3,0){L}
    \tkzDefPoint(3,0){A}

    \tkzDrawCircle(O,R)
    \tkzDrawSegment(L,A)
    \tkzDrawSegment(O,T)
    \tkzDrawPoints(O,T)

    \tkzLabelPoints[above right](O)
    \tkzLabelPoints[below](T)
    \tkzMarkRightAngle[size=0.3](O,T,A)
\end{tikzpicture}
\end{center}

\formula{OT\perp AT.}

\subsection*{Центральный и вписанный углы}

Центральный угол имеет вершину в центре окружности. Его градусная мера равна
градусной мере дуги, на которую он опирается.

Вписанный угол имеет вершину на окружности. Его градусная мера равна половине
градусной меры дуги, на которую он опирается:

\formula{\text{вписанный угол}=\frac12\cdot\text{дуга}.}

Если центральный и вписанный углы опираются на одну и ту же дугу, то

\formula{\text{центральный угол}=2\cdot\text{вписанный угол}.}

Например, центральному углу $60^\circ$ соответствует вписанный угол
$30^\circ$ на той же дуге.

\subsection*{Окружность через три точки}

Через три точки, не лежащие на одной прямой, проходит ровно одна окружность.
Если три точки лежат на одной прямой, окружность через все три точки провести
невозможно.

\section*{\color{darkaccent}13. Стрелки часов}
\sectionline

Пусть время равно $h:m$, где $h$ рассматривается по двенадцатичасовому циклу
($0\le h\le 11$).

Минутная стрелка поворачивается на $6^\circ$ за минуту:

\[
\alpha_{\text{мин}}=6m.
\]

Часовая стрелка поворачивается на $30^\circ$ за час и ещё на $0{,}5^\circ$ за
каждую минуту:

\[
\alpha_{\text{час}}=30h+0{,}5m.
\]

Сначала вычисляем

\[
d=
\left|
\alpha_{\text{мин}}-\alpha_{\text{час}}
\right|.
\]

Меньший угол между стрелками равен

\formula{\min\left(d,360^\circ-d\right).}

Для $3{:}30$:

\[
\alpha_{\text{мин}}=6\cdot30=180^\circ,
\]

\[
\alpha_{\text{час}}
=
30\cdot3+0{,}5\cdot30
=
105^\circ,
\]

\[
d=75^\circ.
\]

\section*{\color{darkaccent}14. Среднее арифметическое, медиана, мода и размах}
\sectionline

Пусть оценки $2$, $3$, $4$, $5$ получили соответственно $1$, $4$, $5$, $2$
учащихся. Упорядоченная выборка имеет вид

\[
2,\;
3,3,3,3,\;
4,4,4,4,4,\;
5,5.
\]

Всего наблюдений $12$.

\subsection*{Среднее арифметическое}

Если значение $x_i$ встречается $n_i$ раз, то

\formula{
\overline{x}
=
\frac{x_1n_1+x_2n_2+\cdots+x_kn_k}
{n_1+n_2+\cdots+n_k}.
}

В нашем примере

\[
\overline{x}
=
\frac{2\cdot1+3\cdot4+4\cdot5+5\cdot2}{12}
=
\frac{44}{12}
=
\frac{11}{3}.
\]

Нельзя вычислять

\[
\frac{2+3+4+5}{4},
\]

поскольку оценки встречаются с разными частотами.

\subsection*{Медиана}

Сначала данные упорядочивают.

Если количество значений нечётное, медиана --- центральное значение. Если
количество значений чётное, медиана --- среднее арифметическое двух
центральных значений.

Здесь центральные позиции --- шестая и седьмая; обе содержат число $4$:

\[
\operatorname{Me}=4.
\]

\subsection*{Мода}

Мода --- значение с наибольшей частотой. В данном примере чаще всего
встречается $4$:

\[
\operatorname{Mo}=4.
\]

\subsection*{Размах}

\formula{R=x_{\max}-x_{\min}.}

\section*{\color{darkaccent}15. Круговые диаграммы}
\sectionline

Если часть величины равна $a$, а вся величина --- $S$, то соответствующий
сектор круговой диаграммы имеет угол

\formula{\alpha=360^\circ\cdot\frac{a}{S}.}

Полезные ориентиры:

\[
\frac14\text{ круга}=90^\circ,
\qquad
\frac12\text{ круга}=180^\circ.
\]

Например, если за шесть дней выставку посетили $12\,000$ человек, то четверть
круга соответствует

\[
\frac{12\,000}{4}=3\,000
\]

посетителей.

\section*{\color{darkaccent}16. Типичные ошибки}
\sectionline

\begin{enumerate}
    \item В неравенстве треугольника проверяется только верхняя граница
    $x<a+b$, а условие $x>|a-b|$ забывается.

    \item При отношении $2:3$ числа $2$ и $3$ перемножаются вместо введения
    величин $2x$ и $3x$.

    \item После нахождения $x$ в задаче на отношения записывается $x$ вместо
    требуемых углов или длин.

    \item Свойство медианы равнобедренного треугольника переносится на
    произвольный треугольник.

    \item Свойство медианы к гипотенузе применяется в непрямоугольном
    треугольнике.

    \item Высоту тупоугольного треугольника пытаются обязательно провести к
    самой стороне, хотя основание высоты может лежать на её продолжении.

    \item Равенство углов определяется только по рисунку без ссылки на
    параллельность, биссектрису или другое свойство.

    \item При вычислении среднего арифметического игнорируются частоты.

    \item Медиана ищется в неупорядоченной выборке.

    \item При избавлении от знаменателя умножается только одно слагаемое
    уравнения вместо всей части уравнения.
\end{enumerate}

\section*{\color{darkaccent}17. Быстрый справочник перед работой}
\sectionline

\begin{center}
\begin{tabularx}{\linewidth}{@{}p{0.38\linewidth}X@{}}
\toprule
\textbf{Ситуация} & \textbf{Что вспоминаем сразу} \\
\midrule
Три стороны треугольника
&
$|a-b|<c<a+b$
\\

Равнобедренный треугольник
&
углы при основании равны; медиана к основанию также высота и биссектриса
\\

Смежные углы
&
сумма $180^\circ$
\\

Биссектрисы смежных углов
&
перпендикулярны
\\

Параллельные прямые
&
соответственные и накрест лежащие равны; односторонние дают $180^\circ$
\\

Прямоугольный треугольник
&
острые углы в сумме $90^\circ$
\\

Угол $30^\circ$
&
противолежащий катет равен половине гипотенузы
\\

Медиана к гипотенузе
&
равна половине гипотенузы
\\

Касательная к окружности
&
радиус в точку касания перпендикулярен касательной
\\

Центральный и вписанный углы
&
на одной дуге центральный угол вдвое больше вписанного
\\

Среднее арифметическое
&
обязательно учитываем частоты
\\

Медиана
&
сначала упорядочиваем данные
\\

Мода
&
значение с наибольшей частотой
\\

Круговая диаграмма
&
$\alpha=360^\circ\cdot\dfrac{\text{часть}}{\text{целое}}$
\\
\bottomrule
\end{tabularx}
\end{center}

\clearpage



\end{document}